\textcolor{red}{Content warning: this paper contains content that may be inappropriate or offensive, such as violence, sexually explicit content, and negative stereotypes and actions.}

\begin{abstract}
Recent advances in diffusion models have significantly enhanced their ability to generate high-quality images and videos, but they have also increased the risk of producing unsafe content.
Existing unlearning/editing-based methods for safe generation remove harmful concepts from models but face several challenges: (1) They cannot instantly remove harmful or undesirable concepts (e.g., artist styles) without additional training. 
(2) Their safe generation capabilities depend on collected training data.
(3) They alter model weights, risking degradation in quality for content unrelated to the targeted toxic concepts.
To address these challenges, we propose \ours{}, a novel, training-free approach for safe text-to-image and video generation, that does not alter the model's weights. 
Specifically, we detect a subspace corresponding to a set of toxic concepts in the text embedding space and steer prompt token embeddings away from this subspace, thereby filtering out harmful content while preserving intended semantics. 
To balance the trade-off between filtering toxicity and preserving safe concepts, \ours{} incorporates a novel self-validating filtering mechanism that dynamically adjusts the denoising steps when applying the filtered embeddings. 
Additionally, we incorporate adaptive re-attention mechanisms within the diffusion latent space to selectively diminish the influence of features related to toxic concepts at the pixel level. 
By integrating filtering across both textual embedding and visual latent spaces, \ours{} ensures coherent safety checking, preserving the fidelity, quality, and safety of the generated outputs.
Empirically, \ours{} achieves \textbf{state-of-the-art} performance in suppressing unsafe content in T2I generation (reducing it by \textbf{22\%} across 5 datasets) compared to other training-free methods and effectively filters targeted concepts, e.g., specific artist styles, while maintaining high-quality output. 
It also shows competitive results against training-based methods. 
We further extend \ours{} to various T2I backbones and T2V tasks, showcasing its flexibility and generalization. 
As generative AI rapidly evolves, \ours{} provides a robust and adaptable safeguard for ensuring safe visual generation.
\end{abstract}
